\section*{Chapter \ref{ch:dv:the-binomial-model} --- The Binomial Model}

\begin{solution}{exo:dv:the-binomial-model:1}
$\Delta=(10-0)/(60-40)=0.5$ share; the deposit $b$ solves $0.5\times40+b=0$, so
$b=-20$ (a loan). Price $0.5\times50-20=5$; $p=(1-0.8)/(1.2-0.8)=0.5$.
\end{solution}

\begin{solution}{exo:dv:the-binomial-model:2}
$u=e^{0.2/\sqrt{12}}=1.0594$, $d=1/u=0.9439$,
$p=(e^{0.03/12}-d)/(u-d)=0.5072$.
\end{solution}

\begin{solution}{exo:dv:the-binomial-model:3}
$C-P=10.3603-4.5925=5.7678$ and $S-KR^{-3}=100-100/1.02^3=5.7678$.
\end{solution}

\begin{solution}{exo:dv:the-binomial-model:4}
The call is built from shares and cash, node by node, for 10.3603, and pays the
call's payoff on every path whatever the probabilities. If it traded above that,
selling it and buying the portfolio would be an arbitrage; below, the reverse.
The colleague's 0.8 changes the call's expected payoff, and equally the share's
expected return: it changes risk premia, not the price relative to the share.
\end{solution}

\begin{solution}{exo:dv:the-binomial-model:5}
It pays at 133.1 and 108.9: $(0.6^3+3\times0.6^2\times0.4)/1.02^3=0.648/1.0612=0.6106$.
\end{solution}

\begin{solution}{exo:dv:the-binomial-model:6}
Cox--Ross--Rubinstein: $u=1.0653$, $d=0.9387$, $p=0.5238$. Jarrow--Rudd:
$u=1.0685$, $d=0.9415$, $p=0.5000$ to four decimals. The first puts the drift in
the probability, the second in the factors.
\end{solution}

\begin{solution}{exo:dv:the-binomial-model:7}
Cox--Ross--Rubinstein 5.5909 (0.0174 too high); Leisen--Reimer 5.5735 (error
$-0.00003$).
\end{solution}

\begin{solution}{exo:dv:the-binomial-model:8}
A probability of 1.125 says that cash grows faster than the share can in either
state ($R=1.06>u=1.05$): shorting the share and depositing the proceeds earns a
riskless profit. The ``calibration'' describes an arbitrage, and prices computed
with it are meaningless. The model requires $d<R<u$.
\end{solution}

\begin{solution}{pb:dv:the-binomial-model:1}
\textbf{1.} $p=(1.02-0.9)/0.2=0.6$.
\textbf{2.} 133.1, 108.9, 89.1, 72.9.
\textbf{3.} 10.3603.
\textbf{4.} 0.6240 share and a deposit of $-52.04$.
\textbf{5.} 4.5925; $C-P=5.7678=100-100/1.02^3$.
\textbf{6.} At 90 after one step (10 against 9.196 held) and at 81 after two (19
against 17.039).
\textbf{7.} 4.9076.
\textbf{8.} 0.3151.
\textbf{9.} No: 10.3603 for both, since exercising a call early gives up the
interest on the strike and the option's remaining value.
\textbf{10.} Exercising pays 10 now; holding is worth 9.196 because the payoff
can at best be 19 two steps later and is often less: the interest on the strike
outweighs the remaining optionality.
\textbf{11.} 5.5735.
\textbf{12.} $-0.0200$ at 100 steps, $+0.0174$ at 101.
\textbf{13.} $-0.00003$.
\textbf{14.} 6.0902; premium 0.5167 over the formula's European value.
\textbf{15.} With the strike at the spot, even $n$ puts a node on the kink and odd
$n$ straddles it: the discretised payoff is alternately too low and too high near
the kink.
\textbf{16.} Cox--Ross--Rubinstein 199 steps; Leisen--Reimer 7 (every odd $n\ge7$
is within one cent).
\textbf{17.} As $n^2$: 2\,001 steps cost about 400 times the work of 100. For
10\,000 options the desk uses the accurate tree at small $n$, or a grid with the
same work shared across strikes.
\textbf{18.} Two possible moves per step, continuous hedging in the limit, no
costs, a known volatility: it prices by replication in a market that is not
complete.
\textbf{19.} 0.3151 on the whiteboard; 199 steps.
\textbf{20.} It is the weight that makes shares and cash reproduce each other's
prices, fixed by $u$, $d$ and $R$, whatever anyone believes about the future.
\end{solution}

\begin{solution}{iq:dv:the-binomial-model:1}
$\Delta=20/40=0.5$ share, borrow 40 (so that $0.5\times80-40=0$): the call costs
$50-40=10$; $p=0.5$.
\iqlookfor{replication, not an expected value under a guessed probability.}
\end{solution}

\begin{solution}{iq:dv:the-binomial-model:2}
Because the option is replicated by shares and cash on every path; its price is
the \omterm{def:dv:no-arbitrage-and-the-fundamental-theorems:claim}{replicating portfolio}'s cost. Any change in the probability changes the
share's expected return as well, and the relative price is unchanged.
\iqlookfor{replication; the probability as a derived weight.}
\end{solution}

\begin{solution}{iq:dv:the-binomial-model:3}
The kink of the payoff sits at a different place relative to the terminal nodes
for even and odd $n$, so the sampled payoff alternates between too high and too
low. Fix: Leisen--Reimer (centre the tree on the strike), average $n$ and $n+1$,
or smooth the payoff over the last step (chapter 22).
\iqlookfor{the kink as the cause.}
\end{solution}

\begin{solution}{iq:dv:the-binomial-model:4}
A deep in-the-money put is worth nearly $K-S$ at expiry; exercising now earns
interest on $K$ and gives up little optionality. For a call, exercising pays $K$
early and gives up the insurance of the optionality, while $C\ge S-Ke^{-r(T-t)}>
S-K$: it is always better sold than exercised.
\iqlookfor{the interest on the strike against the value of optionality.}
\end{solution}

\begin{solution}{iq:dv:the-binomial-model:5}
$(n+1)(n+2)/2$ nodes; backward induction needs only one time slice at a time,
$n+1$ values, so $O(n)$ memory and $O(n^2)$ time.
\iqlookfor{recombination and the one-slice array.}
\end{solution}

\begin{solution}{iq:dv:the-binomial-model:6}
Read them from the nodes after one and two steps: $\Delta\approx(V_1^1-V_1^0)/
(Su-Sd)$ and gamma from the three nodes of step 2, which are centred on $S$ in a
\omterm{def:dv:the-binomial-model:crr}{Cox--Ross--Rubinstein tree} ($Su^2$, $S$, $Sd^2$). Better, start the tree two steps
before today so that the three nodes of today are $Su^2$, $S$, $Sd^2$ and the
differences are centred.
\iqlookfor{Greeks from the tree's own nodes, centred.}
\end{solution}
